\section{Gaussian-localized derivative Laguerre defects}
\label{sec:gaussian-laguerre}

We now return to the field rather than a zero statistic. The following
derivative-indexed first Laguerre defects have a complete source formula,
an all-order acquisition theorem and an event-resolving signal
\cite{GaussianLaguerreProof,LocalizedMetricsAssessment}. They are distinct
from the Gaussian Gram statistic of Section~\ref{sec:localized} and the
one-point measure of Section~\ref{sec:smooth-transverse}. The terminology
does not import an entire-function characterization or a novelty claim;
the classical Laguerre connection is discussed in the existing source
record \cite{Csordas2014PositiveKernels}. The identities below are derived
directly with the normalization \eqref{eq:F-H}.

\subsection{The complete quadratic formula and the kinetic threshold}
Here \(v>0\) always denotes \emph{variance}. For integers \(n\ge0\), set
\begin{equation}\label{eq:gaussian-defect}
 \begin{gathered}
 w_v(y)=\frac{\ee^{-y^2/(2v)}}{\sqrt{2\pi v}},\qquad
 f_n=\partial_x^nH_\lambda,\qquad L_f=(f')^2-ff'',\\
 A_n(\lambda,X,v)=\int_\R w_v(x-X)L_{f_n}(x)\dd x.
 \end{gathered}
\end{equation}
All integrations and differentiations at fixed finite order are justified
by Section~\ref{sec:source}. For real \(s\), write
\begin{equation}\label{eq:laguerre-source-moments}
 M_j(s)=\int_\R |p|^j\ee^{sp^2}\Phi(p)\dd p.
\end{equation}
These complete moments are finite. In the notation \eqref{eq:C-B}--\eqref{eq:J},
for \(L\ge0\), \(\beta>0\) and every integer \(j\ge0\),
\begin{equation}\label{eq:gaussian-moment-order}
 M_j(L)\le j!\beta^{-j}B_0J_{0,2}(L,\beta),\qquad B_0=5C_2(0).
\end{equation}
Indeed \((\beta|p|)^j\le j!\ee^{\beta|p|}\). Thus the complete
radial majorant also exposes the growth in derivative order.

\begin{proposition}[Exact full-source quadratic form]\label{prop:gaussian-source}
For real \(\lambda,X\), \(v>0\) and every \(n\ge0\),
\begin{equation}\label{eq:gaussian-full-source}
 A_n=\frac18\int_{\R^2}(-pq)^n(p-q)^2F_\lambda(p)F_\lambda(q)
       \ee^{\ii X(p+q)-v(p+q)^2/2}\dd p\dd q.
\end{equation}
The phase can equivalently be replaced by \(\cos(X(p+q))\).
In particular the positive-half-line formula is
\begin{align}\label{eq:gaussian-half-source}
 A_0=\frac14\int_0^\infty\!\int_0^\infty F_\lambda(p)F_\lambda(q)
 \bigl[&(p-q)^2\ee^{-v(p+q)^2/2}\cos(X(p+q))\nonumber\\
       +&(p+q)^2\ee^{-v(p-q)^2/2}\cos(X(p-q))\bigr]\dd p\dd q.
\end{align}
For \(A_n\), multiply the first bracket term by \((-1)^n(pq)^n\)
and the second by \((pq)^n\).
\end{proposition}
\begin{proof}
From \(f_n=(\ii^n/2)\int p^nF_\lambda(p)\ee^{\ii xp}\dd p\),
symmetrizing the second product in \((f_n')^2-f_nf_n''\) gives
the multiplier \((-1)^n(pq)^n(p^2+q^2-2pq)/8\).
Gaussian integration gives
\(\int w_v(x-X)\ee^{\ii x(p+q)}\dd x
=\ee^{\ii X(p+q)-v(p+q)^2/2}\).
All products are absolutely integrable. Simultaneous reflection of
\(p,q\) removes the imaginary part, and the four quadrants give
\eqref{eq:gaussian-half-source} with the displayed signs.
\end{proof}

Substitution of the unchanged complete theta source means replacing
\(F_\lambda(p)F_\lambda(q)\), on this positive quadrant, by
\begin{equation}\label{eq:gaussian-theta-product}
 \ee^{\lambda(p^2+q^2)+p+q}
 \sum_{a,b\ge1}P_0(a^2\rho(p))P_0(b^2\rho(q))
                    \ee^{-a^2\rho(p)-b^2\rho(q)}.
\end{equation}
The domain and theta summands are positive; the quadratic formula is
signed. The joint tails below pay both infinite indices and both radial
lines, rather than replacing this source by a finite arithmetic model.

For comparison with field energy, put
\[
 \begin{aligned}
 Z&=\int_\R w_v(x-X)H_\lambda(x)^2\dd x,\\
 E_w&=\int_\R|(H_\lambda\sqrt{w_v(\cdot-X)})'|^2\dd x.
 \end{aligned}
\]
Integration by parts gives the exact identities
\begin{equation}\label{eq:gaussian-kinetic}
 A_0=2E_w-\frac{Z}{2v}
       =\partial_\lambda Z+\tfrac12\partial_{XX}Z
       \qquad\text{(at fixed \(v\)).}
\end{equation}
For details, first
\(A_0=2\int w_vH_x^2-\tfrac12\int w_v''H^2\).
Expanding \(E_w\) and using
\(w_v''/w_v=(x-X)^2/v^2-1/v\) proves the first equality.
The second follows from
\(Z_\lambda=-2\int w_vHH_{xx}\) and
\(Z_{XX}/2=\int w_v(H_x^2+HH_{xx})\).
Consequently \(E_w\ge0\) does not supply the stronger signed threshold
\(E_w\ge Z/(4v)\). The global energy identities already proved remain
regular through the same event.

\subsection{The hierarchy and its complete order series}
\begin{proposition}[Heat and variance identities]\label{prop:gaussian-hierarchy}
For every index in \eqref{eq:gaussian-defect},
\begin{equation}\label{eq:gaussian-hierarchy}
 \partial_\lambda A_n=-\partial_{XX}A_n+2A_{n+1},\qquad
 \partial_v A_n=\tfrac12\partial_{XX}A_n.
\end{equation}
Hence a differentiable positive-variance schedule satisfies
\begin{equation}\label{eq:gaussian-schedule}
 \frac{\dd}{\dd\lambda}A_n(\lambda,X,v(\lambda))
   =(v'/2-1)\partial_{XX}A_n+2A_{n+1}.
\end{equation}
For every \(\epsilon>0\), uniformly in all real \(X\), \(v>0\),
\(\lambda\le L\) and all integers \(n\ge0\),
\begin{equation}\label{eq:gaussian-factorial}
 |A_n|\le n!\epsilon^{-n}\mathcal C_\epsilon(L),\qquad
 \mathcal C_\epsilon(L)=\tfrac14M_0(L+\epsilon/2)M_2(L+\epsilon/2).
\end{equation}
The same bound holds pointwise for \(|L_{f_n}(x)|\).
\end{proposition}
\begin{proof}
For \(f=f_n\), the heat equation gives
\((L_f)_\lambda=(f'')^2+ff''''-2f'f'''\), whereas
\((L_f)''=(f'')^2-ff''''\). Thus
\((L_f)_\lambda=-(L_f)''+2L_{f'}\).
Convolution proves the first identity, and the Gaussian heat equation
proves the second. For the bound use
\[
 |pq|^n\le n!\epsilon^{-n}\ee^{\epsilon|pq|}
        \le n!\epsilon^{-n}\ee^{\epsilon(p^2+q^2)/2}.
\]
For the even positive source
\(V(p)=\ee^{(L+\epsilon/2)p^2}\Phi(p)\),
\(\int(p-q)^2V(p)V(q)\dd p\dd q=2M_0M_2\), because
\(\int pV(p)\dd p=0\). Taking the phase modulus in
\eqref{eq:gaussian-full-source} proves the constant \(1/4\),
also before Gaussian averaging.
\end{proof}

The spatial cancellation characteristic is
\(\lambda=\lambda_0+h\), \(v=v_0+2h>0\). Its exact solution is
\begin{equation}\label{eq:gaussian-characteristic}
 A_n(\lambda_0+h,X,v_0+2h)
   =\sum_{k=0}^\infty\frac{2^kh^k}{k!}A_{n+k}(\lambda_0,X,v_0).
\end{equation}
Indeed the new spectral kernel divided by the old one is
\(\ee^{-2hpq}\). Expand this exponential; absolute domination follows
from \(\ee^{2|h||pq|}\le\ee^{|h|(p^2+q^2)}\).
This proves convergence for both signs of \(h\) in the positive-variance
domain, without iterating a merely finite-order differential identity.
To quantify it, fix \(|h|\le D_h\), \(\lambda_0\le L\),
\(\epsilon>2D_h\), and \(r=2D_h/\epsilon<q_*<1\). For \(K\ge0\),
\begin{align}\label{eq:gaussian-order-tail}
 \left|A_n(\lambda_0+h,X,v_0+2h)
       -\sum_{k=0}^K\frac{2^kh^k}{k!}A_{n+k}(\lambda_0,X,v_0)\right|
 \le n!\epsilon^{-n}\mathcal C_\epsilon(L)
       \frac{(r/q_*)^{K+1}}{(1-q_*)^{n+1}}.
\end{align}
This follows from
\((n+k)!/(n!k!)=\binom{n+k}{k}\), the generating function
\(\sum_{k\ge0}\binom{n+k}{k}q_*^k=(1-q_*)^{-n-1}\), and
\(r^k\le(r/q_*)^{K+1}q_*^k\) for \(k>K\).
The factors depending on \(n\) remain costs in a joint large-order limit.
No uniform positive lower bound on either endpoint variance was used.

\subsection{Joint theta-index, radial and order errors}
The truncation indices satisfy \(N,K\in\mathbb Z_{\ge0}\), and \(U\ge0\).
Use only as an integral approximation the positive reflected cutoff
\begin{equation}\label{eq:gaussian-source-cutoff}
 \Phi_{N,U}(p)=\boldsymbol1_{|p|\le U}\ee^{|p|}
 \sum_{1\le a\le N}P_0(a^2\rho(|p|))\ee^{-a^2\rho(|p|)},
 \qquad0\le\Phi_{N,U}\le\Phi.
\end{equation}
No smooth endpoint extension or operator-domain claim is made about it;
its transform and all finite \(x,\lambda\) derivatives exist. Let
\(A_j^{N,U}\) mean \eqref{eq:gaussian-full-source} with this source.
For \(j=0,2\), \(S\ge0\), set
\[
 \delta_j(S;N,U)=\int_\R|p|^j\ee^{Sp^2}(\Phi-\Phi_{N,U})(p)\dd p,
 \qquad J^{>U}_{j,2}(S,0)=\int_{|p|>U}|p|^j
     \ee^{Sp^2+|p|}\rho(|p|)^2\ee^{-\rho(|p|)}\dd p.
\]
The complete source bounds give
\begin{equation}\label{eq:gaussian-source-tail}
 \delta_j\le5C_2(N)J_{j,2}(S,0)+B_0J^{>U}_{j,2}(S,0),\qquad
 C_2(N)\le\frac{(N+1)^4\ee^{-\pi((N+1)^2-1)}}
                    {1-16\ee^{-\pi(2N+3)}}\quad(N\ge0).
\end{equation}
The index bound is the geometric ratio estimate for consecutive terms;
the denominator is positive. Both errors tend independently to zero.
Writing \(R_U=\pi\ee^{4U}\), the radial payment is explicitly
\begin{equation}\label{eq:gaussian-radial-tail}
 J^{>U}_{j,2}(S,0)=\frac{2}{4^{j+1}}\int_{R_U}^\infty
 \left(\log\frac{s}{\pi}\right)^j
 \left(\frac{s}{\pi}\right)^{1/4}s
 \ee^{S\log^2(s/\pi)/16-s}\dd s.
\end{equation}
For example, when \(R_U\ge\max\{\pi,4S^2/(\ee^4\pi)\}\),
\begin{equation}\label{eq:gaussian-gamma-tail}
 J^{>U}_{j,2}(S,0)\le
 \frac{2c_j}{4^{j+1}\sqrt\pi}\,2^{5/2}\Gamma(5/2,R_U/2),
 \qquad c_0=1,\quad c_2=(8/\ee)^2.
\end{equation}
Here \(\Gamma(5/2,z)=\int_z^\infty u^{3/2}\ee^{-u}\dd u\).
To verify the bound, use \(\log x\le(4/\ee)x^{1/4}\) and
\((\log x)^j\le c_jx^{1/4}\) for \(x\ge1\), \(j=0,2\).
The threshold makes \(S\log^2(s/\pi)/16\le s/2\).

Take \(S=L+\epsilon/2\) and put
\begin{equation}\label{eq:gaussian-T-tail}
 \mathcal T_\epsilon(L;N,U)=\tfrac14
  \{M_2(S)\delta_0(S;N,U)+M_0(S)\delta_2(S;N,U)\}.
\end{equation}
Then \(|A_n-A_n^{N,U}|\le n!\epsilon^{-n}\mathcal T_\epsilon\).
For a proof, the absorbed product difference is nonnegative and even in
each coordinate. Its integral against \((p-q)^2\) equals
\(2\{M_0M_2-M_0^{N,U}M_2^{N,U}\}\), bounded by twice the bracket
in \eqref{eq:gaussian-T-tail}. No cross term is discarded.
Combining this with \eqref{eq:gaussian-order-tail} proves the joint bound
\begin{align}\label{eq:gaussian-joint-tail}
 &\left|A_n(\lambda_0+h,X,v_0+2h)
   -\sum_{k=0}^K\frac{2^kh^k}{k!}A_{n+k}^{N,U}(\lambda_0,X,v_0)\right|
 \nonumber\\
 &\quad\le n!\epsilon^{-n}\left\{
 \mathcal C_\epsilon(L)\frac{(r/q_*)^{K+1}}{(1-q_*)^{n+1}}
 +\frac{\mathcal T_\epsilon(L;N,U)}{(1-r)^{n+1}}\right\}.
\end{align}
The coefficient errors were summed to infinity using the same binomial
identity. Thus every real center, complete index, radial line and order
series is controlled simultaneously. Growing \(n\) requires paying the
displayed factorial and geometric costs by choices of \(K,N,U\);
fixed-order convergence is not a uniform large-order sign theorem.

\subsection{Every finite real multiplicity has a resolving level}
Suppose that at a real event
\begin{equation}\label{eq:gaussian-event}
 \partial_x^jH_{\lambda_*}(x_*)=0\ (0\le j<m),\qquad
 a_m=\partial_x^mH_{\lambda_*}(x_*)\ne0,\qquad m\ge2.
\end{equation}
Fix \(T>0\) such that \([\lambda_*-T,\lambda_*]\) is contained in a
compact real Newman-time interval with upper bound \(L\ge0\).
For \(0\le t\le T\), define
\[
 g(t,y)=\partial_x^{m-2}H_{\lambda_*-t}(x_*+y),\qquad
 C=\tfrac12M_{m+1}(L),\qquad D=\tfrac12M_{m+2}(L).
\]
Thus \(g_t=g_{yy}\), \(g(0,0)=g_y(0,0)=0\), \(g_{yy}(0,0)=a_m\),
and the global real-line bounds are \(|g_{yyy}|\le C\),
\(|g_{yyyy}|\le D\). Equation~\eqref{eq:gaussian-moment-order}
makes the dependence on every finite \(m\) explicit, for example
\(C\le(m+1)!\beta^{-m-1}B_0J_{0,2}(L,\beta)/2\).

\begin{theorem}[Paid finite-multiplicity signal]\label{thm:gaussian-resolution}
For \(0<t\le T\), \(0<\kappa<2\), \(b=|a_m|\) and
\(\mu_1=\sqrt{2\kappa/\pi}\),
\begin{equation}\label{eq:gaussian-resolution}
 A_{m-2}(\lambda_*-t,x_*,\kappa t)
       =a_m^2t(\kappa/2-1)+\mathcal E_m(t),
       \qquad |\mathcal E_m(t)|\le E_\kappa(t;b,C,D),
\end{equation}
where the full error is
\begin{align}\label{eq:gaussian-resolution-error}
 E_\kappa(t;b,C,D)={}&bC\mu_1(4+10\kappa/3)t^{3/2}\nonumber\\
 &+\{bD(3+\kappa)/2+C^2(1+2\kappa+5\kappa^2/4)\}t^2\nonumber\\
 &+CD\mu_1(\kappa/3+3/2)t^{5/2}+D^2t^3/2.
\end{align}
In particular, with \(\delta_\kappa=1-\kappa/2\),
\begin{equation}\label{eq:gaussian-resolution-margin}
 E_\kappa\le\tfrac12\delta_\kappa b^2t
 \quad\Longrightarrow\quad A_{m-2}\le-\tfrac12\delta_\kappa b^2t<0.
\end{equation}
This holds at every fixed finite event for sufficiently small \(t>0\).
\end{theorem}
\begin{proof}
Taylor expansion in space and integration of the heat equation give,
on the \emph{whole} real line,
\begin{align*}
 g&=a_m(y^2/2+t)+R_0,&
 |R_0|&\le C(|y|^3/6+t|y|)+Dt^2/2,\\
 g_y&=a_my+R_1,& |R_1|&\le C(y^2/2+t),\\
 g_{yy}&=a_m+R_2,& |R_2|&\le C|y|+Dt.
\end{align*}
For example, \(|g_{yy}(s,y)-a_m|\le C|y|+Ds\), whose time
integral supplies the first row's time error. The second row uses
\(g_y(t,y)-g_y(0,y)=\int_0^tg_{yyy}(s,y)\dd s\).
The leading defect is \(a_m^2(y^2/2-t)\). With \(y=\sqrt t\,z\),
the absolute remaining terms in \((g_y)^2-gg_{yy}\) are bounded by
\begin{align*}
 &bCt^{3/2}(4|z|+5|z|^3/3)+bDt^2(3/2+z^2/2)\\
 &\quad+C^2t^2(1+2z^2+5z^4/12)
       +CDt^{5/2}(3|z|/2+|z|^3/6)+D^2t^3/2.
\end{align*}
Integrate under \(z\sim N(0,\kappa)\), using
\(Ez^2=\kappa\), \(Ez^4=3\kappa^2\),
\(E|z|=\mu_1\), \(E|z|^3=2\kappa\mu_1\).
This gives every coefficient in \eqref{eq:gaussian-resolution-error}.
The global source bounds pay the Gaussian exterior directly, without
an omitted background or root labeling. Since the error is
\(O(t^{3/2})\) at the fixed event, the last assertion follows.
\end{proof}

The proof applies to any real heat field with the stated global
derivative bounds and finite Gaussian defect, with its own constants;
entire continuation alone is not substituted for integrability.
For the exact heat polynomial
\(H_{\lambda_*-t}(x_*+y)=a_m\ee^{t\partial_y^2}y^m/m!\),
the derivative level is exactly \(a_m(y^2/2+t)\), so the error vanishes
for every \(m\ge2\). For canonical fields, the proof above pays the
complete background. A local-Cauchy version with a separately paid
Gaussian exterior is recorded in \cite{GaussianLaguerreProof}; it is
unneeded for this stronger full-line proof.

After division by \(b^2t\), the error costs are explicitly
\((C/b)\sqrt t\), \(\{D/b+(C/b)^2\}t\),
\((CD/b^2)t^{3/2}\), and \((D/b)^2t^2\), with the displayed
coefficients. Uniform upper source moments give neither a positive
lower bound on \(b\) nor a uniform event radius. No bound on finite
multiplicity was assumed. At fixed positive variance the defect itself
is smooth at the event; shrinking variance and amplitude normalization
are separate costs.

\subsection{Controls, variance conflict and the missing signed estimate}
For the quadratic control \(Q(x)=x^2+d^2\),
\[
 L_Q=2(x^2-d^2),\qquad A_0(Q;X,v)=2(X^2+v-d^2).
\]
Its imaginary pair is detected at \(X=0\) precisely when \(v<d^2\).
The antiheat family \(Q_\lambda=x^2-2(\lambda-\lambda_*)\)
has \(d^2=2t\), \(a_2=2\), and realizes
\eqref{eq:gaussian-resolution} with zero error. For the cubic control
\(C_c(x)=x(x^2+c)\), \(c>0\),
\begin{align*}
 L_{C_c}&=3x^4+c^2>0,& A_0(C_c;0,v)&=9v^2+c^2>0,\\
 L_{C_c'}&=18x^2-6c,& A_1(C_c;X,v)&=18(X^2+v)-6c.
\end{align*}
The first level misses the imaginary pair even without smoothing;
the next level detects it at zero for \(v<c/3\).
The exact family \(C_\lambda=x^3-6(\lambda-\lambda_*)x\) has
\(c=6t\), \(a_3=6\), and \(A_1=36t(\kappa/2-1)\).
These polynomial diagnostics fail the literal complete theta-source
premise and are not proposed canonical fields.

Resolution uses \(v=\kappa(\lambda_*-\lambda)\), hence \(v'=-\kappa\)
and \(\dd A_n/\dd\lambda=-(1+\kappa/2)\partial_{XX}A_n+2A_{n+1}\).
Spatial cancellation instead requires \(v'=2\). A characteristic
ending at the event has \(v=v_*-2t\): for \(v_*>0\) it is outside
\(0<v<2t\) near the event, and for \(v_*=0\) its variance is negative.
Starting at \(t_0>0\), \(v_0=\kappa t_0\), the forward cancellation
characteristic obeys
\begin{equation}\label{eq:gaussian-cone-exit}
 t=t_0-h,\quad v=\kappa t_0+2h,\qquad
 v<2t\ \Longleftrightarrow\ h<(2-\kappa)t_0/4.
\end{equation}
Its exact cone-exit advance is therefore \((2-\kappa)t_0/4\).
Backward continuation has positive variance only for
\(s<\kappa t_0/2\), where \(\lambda=\lambda_0-s\).
The local remainder may further shrink the resolving region.
For \(h\ge0\), the order series conditionally preserves nonnegativity
if \emph{every} required initial \(A_{n+k}\) is nonnegative; those signs
are not supplied by its absolute bound, and resolution is not retained.

The actual unpaid comparison is a lower bound for the signed complete
quadratic form \eqref{eq:gaussian-full-source}. For a relevant event,
one sufficient independently acquired estimate would be
\begin{equation}\label{eq:gaussian-missing-comparison}
 A_{m-2}(\lambda_*-t,x_*,\kappa t)\ge-c|a_m|^2t,
 \quad0\le c<\delta_\kappa,\quad
 E_\kappa(t)<(\delta_\kappa-c)|a_m|^2t.
\end{equation}
It would contradict \eqref{eq:gaussian-resolution}. Neither this
comparison nor unconditional nonnegativity is proved here.
The absolute allowance divided by the desired margin costs
\((m-2)!\epsilon^{-(m-2)}\mathcal C_\epsilon(L)/(|a_m|^2t)\),
which need not be small. This limits that allowance, not the actual
defect or every sharper theta-sensitive estimate.

The collision control of Section~\ref{sec:control} shares regular heat
and energy identities and complete finite moments, yet fails the exact
canonical theta arithmetic. The Euler-factor, weighted-square,
reflected smooth-shape and positive-primitive control classes likewise
share the generic identities when their own complete moment and tail
hypotheses hold; they fail the unchanged canonical coefficient sequence,
integer-square expansion or full reflected endpoint data as applicable.
Variable-kernel controls additionally fail the literal cosine observation
or full heat equation where those were changed. No canonical numerical
constants are assigned to these other classes. Naming their failed
arithmetic premises does not establish
\eqref{eq:gaussian-missing-comparison}. The scoped conclusion excludes an
automatic generic-energy or characteristic argument and leaves
source-sensitive comparisons open.
